% Kapitel 6: Evaluation
% Ergebnisse der SIP-Integration.
\chapter{Evaluation}
\label{chap:evaluation}

In dieser Evaluation werden lediglich Tests aus der \gls{gls-SIP}-Integration behandelt. Eine quantitative Auswertung der Web\gls{gls-UI} selbst ist nicht das Ziel der Evaluation.

\section[SIP-Evaluation]{\gls{gls-SIP}-Evaluation}
Die \gls{gls-SIP}-Anbindung wurde im Gegensatz zur Web\gls{gls-UI} detailliert verfolgt und ausgewertet. Die Signalisierungsebene und der Medienpfad funktionieren stabil: \gls{gls-SIP}-Trunk-Registration beim Provider, Zuteilungsregeln und die automatische Raumerstellung greifen zuverlässig. Eingehende Anrufe werden korrekt erkannt und an den \gls{gls-KI}-Agenten weitergeleitet.

Nach den durchgeführten \gls{gls-SIP}-Testanrufen zeigt sich, dass:

\begin{itemize}
  \item Bidirektionale Audioübertragung ohne \gls{gls-RTP}-Timeout erfolgt
  \item Mehrfache Gesprächsumläufe stabil verarbeitet (keine Abbrüche oder Timeouts)
  \item Anrufer Begrü\ss ung und Sprachausgabe des Agents deutlich empfangen
  \item Audio-Qualität und Latenz im erwarteten Bereich liegen
\end{itemize}

\subsection{Integration, Durchführung und Testszenarien}

Die \gls{gls-SIP}-Integration wurde in einer kontrollierten \gls{gls-LAN}-Umgebung mit Docker Desktop betrieben, während die Testanrufe selbst über das öffentliche Mobilfunknetz (Telekom) durchgeführt wurden. Aus der Gesamtheit der Testanrufe wurden exemplarisch 5 Anrufe für eine detaillierte Analyse ausgewählt. Die folgenden Beispiele zeigen die Vielfalt der getesteten Szenarien:

\begin{itemize}
  \item \textbf{Anruf 1:} Fragen zu physikalischen Dimensionen (Fu\ss ballfeld, Handball-Umfang): 3 Turns, 54,97s Dauer
  \item \textbf{Anruf 2:} Geographie und Zeitberechnungen (tiefste Meeresstelle): 3 Turns, 56,75s Dauer
  \item \textbf{Anruf 3:} Infrastruktur-Fragen (Deutsche Autobahnen): 3 Turns, 49,0s Dauer
  \item \textbf{Anruf 4:} Tier-Rekorde (schnellstes/langsamstes Tier): 4 Turns, 46,96s Dauer
  \item \textbf{Anruf 5:} Technologie-Abkürzungen (IT-Bereiche): 3 Turns, 37,96s Dauer
\end{itemize}

Die Testanrufe wurden bewusst mit unterschiedlicher Komplexität und Anzahl an Turns gestaltet , um eine Varianz in der Ergebnisdarstellung zu erreichen. Alle Anrufe wurden regulär durch Auflegen des Anrufers beendet (CLIENT\_INITIATED disconnect).

\subsection{Metriken und Analyse}

Aus den durchgeführten Testanrufen wurden fünf Anrufe für eine detaillierte Metrikanalyse ausgewählt. Sitzungen ohne Turns (z.\,B. kurz abgebrochene Rufaufbauten durch die Testperson) wurden von der Aggregation ausgeschlossen.

Zur qualitativen Einordnung der \gls{gls-STT}-Robustheit werden exemplarisch die in den Logs erkannten Nutzerfragen der fünf Testanrufe aufgeführt; dabei ist eine einzelne Fehlinterpretation sichtbar, die die Grenzen der Erkennung in Telefonszenarien verdeutlicht:
\begin{itemize}
  \item \textbf{Anruf 1 (Physikalische Dimensionen):} „Hallo, wie groß ist ein Fu\ss ballfeld?“, „Gib mir den Umfang für einen Handball.“, „Danke.“
  \item \textbf{Anruf 2 (Geographie/Zeiten):} „Gib mir die tiefste Stelle im Meer.“, „Wie lange würde es dauern, dort hinunterzutauchen?“, „Danke.“
  \item \textbf{Anruf 3 (Infrastruktur):} „Wie lang ist die längste Autobahn in Deutschland?“, „Welche ist die kürzeste?“, „Danke.“
  \item \textbf{Anruf 4 (Tier-Rekorde):} „Welches ist das schnellste Tier der Welt?“, \textit{„Welche Systemstelle?“} (Fehlinterpretation), „Das Langsamste.“, „Danke.“
  \item \textbf{Anruf 5 (Technologie-Abkürzungen):} „Was ist die Abkürzung IT?“, „Gib mir einen Bereich davon.“, „Danke.“
\end{itemize}

Die folgende Tabelle~\ref{tab:sip-metriken-aggregiert} zeigt die aggregierten Messwerte über diese fünf exemplarischen Testanrufe hinweg:

\begin{table}[H]
  \centering
  \begin{tabular}{|l|c|c|c|c|}
    \hline
    \textbf{Metrik} & \textbf{Mittelwert} & \textbf{Median} & \textbf{Min} & \textbf{Max} \\ \hline
    Anruf-Dauer (in s) & 49,13 & 49,00 & 37,96 & 56,75 \\ \hline
    Turns pro Anruf & 3,2 & 3 & 3 & 4 \\ \hline
    STT-Latenz (Ø in s) & 0,86 & 0,78 & 0,52 & 1,19 \\ \hline
    Transcript-Delay (Ø in s) & 1,30 & 1,31 & 1,05 & 1,74 \\ \hline
    TTS-Latenz (Ø in s) & 1,45 & 1,41 & 1,15 & 1,83 \\ \hline
    Gesamt Audio (s) & 21,15 & 22,35 & 10,73 & 26,78 \\ \hline
    Audio-Bytes (Ø in KiB) & 495,6 & 523,7 & 251,6 & 627,6 \\ \hline
  \end{tabular}
  \caption{Aggregierte Metriken und Latenzen der \gls{gls-SIP}-Tests}
  \label{tab:sip-metriken-aggregiert}
\end{table}

% --- SIP-Aggregat: Kurzverweis auf Statistik, Fokus auf Netzwerk- und Protokollaspekte ---
Die Analyse der aggregierten Daten zeigt konsistente Leistungswerte über die fünf detailliert ausgewerteten Testanrufe hinweg. Die \gls{gls-STT}-Latenz variiert zwischen 0,52 s im besten Fall und 1,19 s bei längeren Prompts. Im \gls{gls-SIP}-Kontext treten zusätzliche Latenzen durch Netzwerk- und Protokollaspekte auf, insbesondere bei der Transkriptverzögerung, die mit durchschnittlich 1,30 s deutlich über der reinen \gls{gls-STT}-Latenz liegt \parencite{ITU-T-G114}.

Die \gls{gls-TTS}-Latenz bleibt mit durchschnittlich 1,45 s über alle Anrufe hinweg konsistent \parencite{openai_whisper_docs}. Die durchschnittlich verarbeiteten Audiodaten pro Anruf liegen mit rund 496 KiB im niedrigen Kilobyte-Bereich. Dies entspricht den Empfehlungen für effiziente Sprachübertragung mit modernen Codecs wie Opus \parencite{rfc6716}.

\subsection{Detailanalyse von Testanrufen}

Zur tieferen Einordnung werden im Folgenden die Kennzahlen der fünf exemplarisch ausgewählten Testanrufe dargestellt (Tabelle~\ref{tab:sip-calls-detail}):

\begin{table}[H]
  \centering
  \small
  \begin{tabular}{|l|c|c|c|c|c|}
    \hline
    \textbf{Metrik} & \textbf{Anruf 1} & \textbf{Anruf 2} & \textbf{Anruf 3} & \textbf{Anruf 4} & \textbf{Anruf 5} \\ \hline
    Dauer (in s) & 54,97 & 56,75 & 49,00 & 46,96 & 37,96 \\ \hline
    Turns (Anzahl) & 3 & 3 & 3 & 4 & 3 \\ \hline
    STT-Durchschnitt (in s) & 1,19 & 0,77 & 0,52 & 1,07 & 0,78 \\ \hline
    TTS-Durchschnitt (in s) & 1,83 & 1,58 & 1,29 & 1,15 & 1,41 \\ \hline
    Audio (in s) & 25,07 & 26,78 & 22,35 & 20,81 & 10,73 \\ \hline
    Bytes (in KiB) & 587,6 & 627,6 & 523,7 & 487,7 & 251,6 \\ \hline
    STT-Fehlinterpretationen (Anzahl) & 0 & 0 & 0 & 1 & 0 \\ \hline
  \end{tabular}
  \caption{Detaillierte Metriken der \gls{gls-SIP}-Tests}
  \label{tab:sip-calls-detail}
\end{table}

% --- Detailanalyse: Fokus auf Gesprächsdynamik und Nutzerverhalten ---
Die Detailanalyse der einzelnen Anrufe zeigt Unterschiede in den gemessenen Leistungsmetriken. In Anruf 4 trat eine einzelne Fehlinterpretation des Transkripts auf, während in den übrigen Anrufen keine weiteren Auffälligkeiten dokumentiert wurden. Anruf 1 und Anruf 4 weisen höhere \gls{gls-STT}-Latenzen auf (1,19 s bzw. 1,07 s), während Anruf 3 mit 0,52 s die geringste \gls{gls-STT}-Latenz erreicht. Die \gls{gls-TTS}-Latenz variiert zwischen 1,15 s und 1,83 s.

Anruf 5 weist mit einer Gesamtdauer von 37,96 s sowie der geringsten Audiodatenmenge (10,73 s und 257,6 KiB) die niedrigsten Werte auf. Die beobachteten Unterschiede verdeutlichen die Variabilität der Systemmetriken zwischen einzelnen Gesprächen, deren Länge und der Anzahl Turns \parencite{Data2025_AIDataQuality,ssrn5348245}.

\subsection{Dialogbasierte Latenzanalyse}

Um die Konsistenz der Leistung innerhalb einzelner Gespräche zu bewerten, wurde eine dialogbasierte Analyse durchgeführt (Tabelle~\ref{tab:turn-latencies}):

\begin{table}[H]
  \centering
  \begin{tabular}{|l|c|c|c|c|}
    \hline
    \textbf{Anruf} & \textbf{Turn 1 STT (s)} & \textbf{Turn 2 STT (s)} & \textbf{Turn 3 STT (s)} \\ \hline
    Anruf 1 & 1,02 & 2,10 & 0,46  \\ \hline
    Anruf 2 & 0,84 & 0,91 & 0,55  \\ \hline
    Anruf 3 & 0,82 & 0,38 & 0,35  \\ \hline
    Anruf 4 & 1,07 & 0,81 & 1,80  \\ \hline
    Anruf 5 & 0,68 & 1,10 & 0,54  \\ \hline
  \end{tabular}
  \caption{Dialogbasierte Latenzanalyse einzelner \gls{gls-SIP}-Tests}
  \label{tab:turn-latencies}
\end{table}

% --- dialogbasiert: Fokus auf Schwankungen und Robustheit, keine Wiederholung der Statistik-Argumentation ---
Die dialogbasierte Analyse offenbart eine erhebliche Varianz der \gls{gls-STT}-Latenz innerhalb einzelner Anrufe. In Anruf 1 steigt die Latenz im zweiten Turn auf 2,10 s an, sinkt jedoch im dritten Turn wieder auf 0,46 s ab. Anruf 4 weist mit 1,80 s im dritten Turn eine besonders hohe Varianz auf. Im Gegensatz dazu bleibt Anruf 3 durchgehend konsistent und unterschreitet in allen Turns die 0,85 s. Die Ratio zwischen dem schnellsten und dem langsamsten Turn liegt bei 6,1:1 (0,35 s vs 2,10 s), was im erwartbaren Bereich für cloudbasierte \gls{gls-KI}-Dienste liegt \parencite{whisper_quant_2025,ITU-T-G114}.

Alle Turns wurden erfolgreich verarbeitet, ohne dass Timeouts auftraten. Eine Fehlinterpretation wurde in Anruf 4 dokumentiert.

\subsection{Audioqualität und Medienpfad-Stabilität}

Die Stabilität des Medienpfads wurde durch kontinuierliche bidirektionale Audioübertragung über die gesamte Anrufdauer nachgewiesen. Folgende qualitative Beobachtungen wurden während aller \gls{gls-SIP}-Testanrufe dokumentiert:

\begin{itemize}
  \item \textbf{Codec:} Opus 48 kHz, 24 kbps \parencite{rfc6716}
  \item \textbf{Paketverlust:} In \gls{gls-RTP}-Statistiken/LiveKit-Logs keine Paketverluste beobachtet
  \item \textbf{Jitter:} Stabile Audioausgabe ohne starke Verzögerungen
  \item \textbf{Echo:} Keine Echo-Bildung beobachtet (subjektive Wahrnehmung)
  \item \textbf{Verständlichkeit:} Subjektiv geschätzte \gls{gls-MOS} Werte von 4,0--4,5 / 5,0 (gut bis sehr gut)
  \item \textbf{\gls{gls-RTP}-Timeout:} Kein Auftreten bei allen Testanrufe
\end{itemize}

% --- Audioqualität: Fokus auf subjektive und objektive Bewertung, Codec, Netzwerkstabilität ---
Die konsistente Audioqualität über alle Testanrufe hinweg bestätigt die Stabilität des Medienpfads. Die \gls{gls-MOS}-Einschätzung ist subjektiv und basiert nicht auf standardisierten ITU-T-Hörversuchen. Neben der technischen Stabilität ist insbesondere die subjektive Verständlichkeit ein zentrales Qualitätsmerkmal für den Praxiseinsatz \parencite{ITU-T-G114}. Die Nutzung des Opus-Codecs und das Fehlen von Paketverlusten oder Jitter unterstreichen die objektive und subjektive Güte der Audioübertragung \parencite{rfc6716}.

\section{Hypothesenbezogene Evaluationsergebnisse}
\label{sec:hypothesen-resultate}

Die erhobenen Daten ermöglichen eine direkte Prüfung der in Kapitel~\ref{chap:einleitung} formulierten Hypothesen. Im Folgenden werden die Ergebnisse kurz dokumentiert, ohne inhaltliche Interpretation:

\paragraph{H1 (Funktionalität):} Gestützt durch Daten.
\begin{itemize}
  \item Keine funktionellen Fehler oder \gls{gls-RTP}-Timeouts über alle Testanrufe
  \item Eine \gls{gls-STT}-Fehlinterpretation wurde dokumentiert (Anruf 4)
  \item Bidirektionale Audioübertragung ohne Medienpfad-Fehler
  \item Alle Verbdindungsabbrüche regulär durch Auflegen des Anrufers (CLIENT\_INITIATED)
\end{itemize}

\paragraph{H2 (Latenz):} Gestützt durch Daten.
\begin{itemize}
  \item \gls{gls-STT}-Latenz: Ø 0,86\,s (Median: 0,78\,s, Spanne: 0,52--1,19\,s)
  \item \gls{gls-TTS}-Latenz: Ø 1,45\,s (Median 1,41\,s, Spanne: 1,15--1,83\,s)
  \item Transkriptverzögerung: Ø 1,30\,s
\end{itemize}

\paragraph{H3 (Ressourceneffizienz):} Gestützt durch Daten.
\begin{itemize}
  \item \gls{gls-CPU}-Auslastung: Ø 11{,}2\,\% (Max. 26\,\%)
  \item\gls{gls-RAM}-Nutzung: 11{,}2--12{,}1\,GB bei verfügbaren 16\,GB
\end{itemize}

\paragraph{H4 (Skalierbarkeit):} Gestützt durch Architekturdesign.
\begin{itemize}
  \item Modulare Komponentenstruktur
  \item Flexible Ressourcenzuteilung durch Containerbasierung
  \item Reserve in \gls{gls-CPU} und \gls{gls-RAM} für Erweiterbarkeit
\end{itemize}

\paragraph{H5 (STT-Robustheit):} Qualitativ gestützt durch Daten.
\begin{itemize}
  \item Die Mehrheit der deutschsprachigen Testprompts wurde korrekt erkannt
  \item Eine einzelne Fehlinterpretation wurde dokumentiert (Anruf 4)
  \item Es wurden keine systematischen Transkriptionsfehler dokumentiert
  \item Eine quantitative Messung der Word-Error-Rate wurde nicht durchgeführt
\end{itemize}

\section{Bedrohungen der Validität}
\label{sec:threats-to-validity}

Die Validität der Evaluationsergebnisse unterliegt mehreren Einschränkungen, die bei der Interpretation der Daten berücksichtigt werden müssen.

\subsection{Interne Validität}
\textbf{Stichprobengrö\ss e:}\\
Die Evaluation basiert auf \gls{gls-SIP}-Testanrufen, wovon fünf Anrufe detailliert analysiert wurden. Für produktive Systeme wären grö\ss ere Testszenarien erforderlich, um statistische Signifikanz zu gewährleisten. Dies ist für das hier getestete prototypische System nicht umgesetzt.

\textbf{Testszenario-Diversität:}\\
Die Tests wurden in einem mit Docker betriebenen System in einer kontrollierten \gls{gls-LAN}-Umgebung durchgeführt, wobei die \gls{gls-SIP}-Anrufe über eine Mobilfunk-Verbindung initiiert wurden. Andere Architekturen können abweichende Ergebnisse liefern.

\textbf{Prompt-Komplexität:}\\
Die verwendeten Testprompts decken Themen des Allgemeinwissen ab (Geographie,
Physik, Technologie). Domänenspezifische Fachsprache oder stark akzentuierte Aussprache könnte die \gls{gls-STT}-Genauigkeit beeinflussen.

\subsection{Externe Validität}
\textbf{Hardware-Abhängigkeit:}\\
Die Messungen wurden auf einem spezifischen Testsystem (16\,GB \gls{gls-RAM},
8 Kerne) durchgeführt. Die Übertragbarkeit auf andere Hardware-Konfigurationen (z.\,B. ARM-Architekturen, Cloud-Instanzen) ist nicht geprüft.

\textbf{Cloud-\gls{gls-API}-Variabilität:}\\
Die gemessenen Latenzen für \gls{gls-STT} und
\gls{gls-TTS} hängen von der Auslastung der OpenAI-Server und der geografischen Distanz ab. Wiederholte Messungen zu anderen Zeitpunkten oder von anderen Standorten könnten abweichende Werte liefern.

\textbf{Telefonanbieter-Spezifika:}\\
Die \gls{gls-SIP}-Tests erfolgten mit einem spezifischen Provider und einem dazugehörigen \gls{gls-SIP}-Trunk. Andere Provider könnten abweichende Codec-Unterstützung oder Netzwerkcharakteristika aufweisen.

\subsection{Konstrukt-Validität}
\textbf{Subjektive \gls{gls-MOS}-Bewertung:}\\
Die \gls{gls-MOS}-Werte basieren auf subjektiver Einschätzung ohne standardisiertes \gls{gls-ITU-T}-Testprotokoll. Eine objektive Messung mittels \gls{gls-PESQ} oder \gls{gls-POLQA} ist nicht Bestandteil der Hypothesen dieser Arbeit.

\textbf{Word-Error-Rate-Messung:}\\
Die qualitative Beobachtung, dass alle Prompts korrekt erkannt wurden, ersetzt keine quantitative Messung der Word-Error-Rate mit Referenztranskripten. Dies ist kein Bestandteil der Hypothesen dieser Arbeit.

\subsection{Reliabilität}
\textbf{Reproduzierbarkeit:}\\
Alle Tests sind durch die bereitgestellten Komponenten und deren Konfigurationen reproduzierbar. Die Cloud-\gls{gls-API}-Abhängigkeit führt jedoch zu nicht-deterministischen Latenzen bei Wiederholungen.

\textbf{Logging-Präzision:}\\
Die Messungen der Latenzen basieren auf Log-Zeitstempeln mit Millisekunden-Präzision. Systematische Logging-Verzögerungen, die durch Region oder Systemabhängigkeiten entstehen, könnten die Messungen verfälschen.